\notepage{72}{Selection rules}
\smallheading{Parity}
For a rotationally invariant $H^0$, label states by $\ket{\alpha\ell m}$.
Their angular dependence is
\[
 \braket{\mathbf r}{\alpha\ell m}\propto Y_\ell^m(\theta,\phi),
 \qquad Y_\ell^m(\pi-\theta,\pi+\phi)=(-1)^\ell Y_\ell^m(\theta,\phi).
\]
Thus parity acts according to
\[
 \mel{\mathbf r}{\Pi}{\alpha\ell m}
   =(-1)^\ell\braket{\mathbf r}{\alpha\ell m}.
\]
Decompose a perturbation into its even and odd parts,
\[
 H'=H'_0+H'_1,\qquad
 \Pi H'_p\Pi^{-1}=(-1)^pH'_p,\qquad p=0,1.
\]
Then
\[
 \mel{\alpha_2\ell_2m_2}{H'_p}{\alpha_1\ell_1m_1}=0
 \quad\text{if }p+\ell_1+\ell_2\text{ is odd}.
\]
If the sum is even, parity alone does not determine whether the matrix
element vanishes.
\smallheading{Angular momentum and spherical tensors}
A perturbation may also be expanded into spherical-tensor components
$T_L^M$. A rank-$L$ tensor component carries the angular-momentum labels
$(L,M)$. For a multiplicative angular tensor, its coordinate action has
the form
\[
 \mel{\theta,\phi}{T_L^M}{\psi}
 \propto Y_L^M(\theta,\phi)\braket{\theta,\phi}{\psi}.
\]
The angular-momentum selection rules following from the Wigner--Eckart
theorem are
\[
 \mel{\alpha_2\ell_2m_2}{T_L^M}{\alpha_1\ell_1m_1}=0
 \quad\text{unless}\quad
 \begin{cases}
 |\ell_1-L|\le\ell_2\le\ell_1+L,\\
 m_2=m_1+M.
 \end{cases}
\]
\editorial{The coordinate multiplication formula describes a
multiplicative tensor, such as an angular factor in a potential; it is
not the definition of every spherical-tensor operator. The parity label
$p$ and the tensor rank $L$ are kept distinct here.}
